\documentclass[11pt,a4paper]{article}
\usepackage[utf8]{inputenc}
\usepackage[T1]{fontenc}
\usepackage[spanish]{babel}
\usepackage{hyperref}
\usepackage{geometry}
\geometry{margin=2.5cm}
\usepackage{parskip}

\title{109 un jugo publicado sin decir quien lo tuvo parece de nadie}
\author{Alejandro Valenzuela Rivera\\\small Universidad Autónoma de Nuevo León}
\date{\today}
\begin{document}
\maketitle

\begin{quote}
\textbf{Sesgo declarado.} Este trabajo se sostiene desde el sesgo declarado de su autor, y no finge neutralidad. Lo que afirma se puede tumbar, y ese es el punto: el falsificador está escrito, no escondido.
\end{quote}

\section*{Resumen}
Cada jugo que se publique lleva al inicio su origen, su autoria, su licencia y su limpieza. Y no es adorno: es lo que lo vuelve citable y lo que impide que se lea como material huerfano.

\section{Por qué se sostiene}
Y es la misma regla que el repo aplica al revés: no se usa un nombre ajeno como propio, y no se publica un texto propio como ajeno. El bloque hace las dos cosas: dice de donde viene el trabajo y con que licencia se puede usar.

\section{De dónde viene}
La autoridad, y el ATTRIBUTION del repo.

\section{El caso}
Ninguna: es la regla de la publicacion.

\section{Método}
Un jugo publicado con su autor, y con lo que pregunto.

\section{Límites, y qué lo tumba}
Una idea publicada sin decir de quien salio.

\section{Lo que queda abierto}
\textquestiondown{}Que falta cuando un texto publicado no dice de quien es?

\section*{Cómo se leyó esta tesis}
Esto se escribió en el repositorio man-made-empathy, del que este candado es una pieza: candados/03 el nombre y la obra/109 un jugo publicado sin decir quien lo tuvo parece de nadie.md. La máquina que lo compuso lo hizo desde este mismo archivo, y no agregó nada que no estuviera escrito. Si algo de esto es falso, se dice en el mismo sitio, y se corrige con fecha.

\end{document}
